\section{Group generators and augmentation ideals}
\label{sec:algebra}

We need finite groups whose group generators are numerous but whose
augmentation ideals have uniformly short generating lists.  The distinction
between these two numbers is classical; see
\cite[Sections~1--2]{GruenbergLinnell}.  We give the estimates needed here,
then prove the elementary row reduction used in the handle construction.
All modules in this section are right modules.

For a finitely generated group $G$, let $d(G)$ be its minimum number of
group generators.

\begin{lemma}\label{lem:group-gap}
Let $S$ be a nontrivial finite group, let $k$ be a finite positive
integer, and let $m\ge0$ be an integer.
If $k>|S|^m$, then $d(S^k)>m$.
\end{lemma}

\begin{proof}
Suppose that $g_1,\ldots,g_m$ generate $S^k$, padding a shorter list with
identities.  Each factor coordinate $i$ gives a tuple
$((g_1)_i,\ldots,(g_m)_i)\in S^m$.  Since $k>|S|^m$, two coordinates
$i\ne j$ have the same tuple.  Every listed generator then belongs to the
proper subgroup
\[
  \{g\in S^k:g_i=g_j\},
\]
a contradiction.
\end{proof}

We shall use
\begin{equation}\label{eq:group-choice}
  b=2,\qquad S=A_5,\qquad G=S^k.
\end{equation}
For now $k\ge1$ is arbitrary.  The group $S$ is nonabelian simple, hence
perfect, and the direct product $G$ is also perfect.  Every finite group is finitely presented: one may
take symbols $t_g$ for its elements and relations $t_1=1$ and
$t_gt_h=t_{gh}$.  Thus these groups satisfy the finite-presentation
condition needed for the symplectic realization below.

Write $\Lambda=\Z[G]$ and let
\[
  \operatorname{aug}:\Lambda\longrightarrow\Z,
  \qquad \sum_{g\in G}a_gg\longmapsto\sum_{g\in G}a_g
\]
be the augmentation map.  Its kernel $I$ is the augmentation ideal,
regarded as a right $\Lambda$-module.

The uniform bound below is a special case of the direct-product theorem
of Cossey, Gruenberg and Kov\'acs~\cite[Theorem~2]{CGK1974}.
For finite groups, that theorem expresses the augmentation-ideal
generator number of a direct product as the maximum of the corresponding
factor numbers and the generator number of its abelianization.
We give an elementary proof for the present family, including the
passage from modular to integral generators.

\begin{lemma}\label{lem:augmentation}
For every finite positive integer $k$, the right augmentation ideal of
$\Z[A_5^k]$ can be generated by two elements.
\end{lemma}

\begin{proof}
We first prove that two generators suffice modulo every prime.
The group $A_5$ is generated by
$a=(1\,2\,3\,4\,5)$ and $c=(1\,2\,3)$.
Indeed, $aca^{-1}=(2\,3\,4)$, and the subgroup generated by
$c$ and $aca^{-1}$ is transitive on $\{1,2,3,4\}$, and the
stabilizer of $4$ contains the order-three group generated by $c$. It is therefore all of the copy
of $A_4$ fixing $5$. The subgroup generated by $a,c$ is transitive
on five letters and has a stabilizer containing this $A_4$, so it
has order at least $5\cdot12=60$ and equals $A_5$.

Put $\Lambda=\Z[G]$ and $I=\ker(\aug:\Lambda\to\Z)$.
The splitting $\Lambda=I\oplus\Z\cdot1$ as abelian groups identifies
\[
 R_p=\F_p[G],\qquad I_p=I/pI=\ker(R_p\xrightarrow{\aug}\F_p).
\]
For each simple right $R_p$-module $V$, we claim
\begin{equation}\label{eq:augmentation-hom-bound}
 \dim_{\F_p}\Hom_{R_p}(I_p,V)\le 2\dim_{\F_p}V.
\end{equation}
If $G$ acts trivially on $V$, every homomorphism $\ell:I_p\to V$
defines a group homomorphism $g\mapsto\ell(g-1)$ to the additive
group of $V$. Since $G$ is perfect, this map vanishes, and thus
$\ell=0$.

Otherwise some factor $S_i\cong A_5$ acts nontrivially on $V$.
Normality of $S_i$ makes $V^{S_i}$ a submodule, so simplicity gives
$V^{S_i}=0$. Write $a_i,c_i$ for the two displayed generators in this
factor. We show that $\ell$ is determined by
$\ell(a_i-1)$ and $\ell(c_i-1)$. The identity
\[
 \ell(gh-1)=\ell(g-1)h+\ell(h-1)
\]
shows that if both values vanish, then $\ell(s-1)=0$ for every
$s\in S_i$. If $h$ lies in the product of the other factors, then
for every $s\in S_i$,
\[
 \ell(h-1)(s-1)
 =\ell((h-1)(s-1))
 =\ell((s-1)(h-1))=0.
\]
Thus $\ell(h-1)\in V^{S_i}=0$. Every element of $G$ has the form
$sh$, so the same identity gives $\ell=0$. Evaluation at the two
generators therefore embeds $\Hom_{R_p}(I_p,V)$ into $V^2$, proving
\eqref{eq:augmentation-hom-bound}.

Let $J=\Jac(R_p)$. We use the elementary finite-dimensional algebra
facts that $J$ annihilates every simple module, $J$ is nilpotent,
and $R_p/J$ is semisimple. For completeness, $J$ is the intersection
of the maximal right ideals, equivalently the annihilators of the
simple right modules. Multiplication by $J$ lowers a composition
series of the regular module, proving nilpotence. Finite dimension
allows finitely many maximal right ideals with intersection $J$;
thus $R_p/J$ embeds in a finite direct sum of simple quotients and is
semisimple, as are its finite-dimensional modules.

Let $a_V,t_V$ denote the multiplicities of a simple module $V$ in
$I_p/I_pJ$ and $R_p/J$, respectively, and let
$e_V=\dim_{\F_p}\End_{R_p}(V)$. Schur's lemma and evaluation at $1$
give
\[
 \dim_{\F_p}\Hom_{R_p}(I_p,V)=a_Ve_V,
 \qquad
 \dim_{\F_p}\Hom_{R_p}(R_p,V)=t_Ve_V=\dim_{\F_p}V.
\]
Consequently $a_V\le2t_V$ for every $V$. There is a surjection
$(R_p/J)^2\to I_p/I_pJ$. Lift the two generator images to $I_p$
and let $N_p$ be their right span. Then
$I_p/N_p=(I_p/N_p)J$, and nilpotence gives $N_p=I_p$.
Hence $I_p$ has a generating pair for every prime $p$.

It remains to lift to two integral generators. We first choose one
element whose right ideal has finite additive index in $I$, while
retaining a residue that extends to a generating pair at each prime
dividing $|G|$. A second element will then remove the finite quotient.
Put $n=|G|=60^k$.
For each $p\in\{2,3,5\}$ choose a generating pair
$(u_p,v_p)$ of $I_p$. By the Chinese remainder theorem on the free
abelian group $I$, choose $y_0\in I$ with residue $u_p$ at each of
these three primes. Define
\[
 z=n\cdot1-\sum_{g\in G}g\in I.
\]
For $x\in I$, we have $zx=nx$. Thus left multiplication on $I$ by
$y_0+30t z$ has matrix $L_{y_0}+30tn\,\mathrm{id}$. Its determinant
is a nonzero polynomial in $t$, so there is an integer $t$ for which
it is nonzero. Fix this $t$ and put
\[
 y_1=y_0+30t z,\qquad D=|\det(L_{y_1}:I\to I)|>0.
\]
The adjugate identity gives $DI\subset y_1I\subset y_1\Lambda$.
Moreover, $y_1$ still reduces to $u_p$ for $p=2,3,5$.

For each prime $p\mid D$, the right module $I_p/y_1R_p$ is cyclic.
For $p\mid n$, it is generated by the image of $v_p$.
For $p\nmid n$, the element
\[
 e_p=1-n^{-1}\sum_{g\in G}g\in R_p
\]
satisfies $I_p=e_pR_p$, so its image generates every quotient of
$I_p$, including $I_p/y_1R_p$. Apply the Chinese remainder theorem on $I$
to choose $y_2\in I$ reducing to $v_p$ when $p\mid D$ and $p\mid n$,
and to $e_p$ when $p\mid D$ and $p\nmid n$. When $D=1$,
take $y_2=0$.

The quotient $C=I/(y_1\Lambda+y_2\Lambda)$ is a finite abelian group
killed by $D$, and the residue choices give $C=pC$ for every prime
$p\mid D$. Iterating these equalities through the prime factorization
of $D$ yields $C=DC=0$. Therefore $y_1,y_2$ generate $I$ as a right
$\Lambda$-module.
\end{proof}

In particular, choosing any finite $k>60^{16b}$ gives
\begin{equation}\label{eq:generator-gap}
  d(G)>16b,
  \qquad I\text{ has a generating list of length }b.
\end{equation}
We always pad such a list by zeros if necessary.  Increasing $k$ makes
$d(G)$ unbounded by Lemma~\ref{lem:group-gap}, while
Lemma~\ref{lem:augmentation} keeps the same bound $b$.

\subsection{An elementary row reduction}

The next lemma supplies an entry equal to $1$ using only elementary
basis additions.  This is the form needed for handle slides.  It is a
special case of Bass's stable-range results applied to the opposite
ring~\cite[Lemma~6.4 and Theorems~4.2(a),~11.1]{Bass}; we include a direct
proof to retain control of the multiplication sides and the row length.

\begin{lemma}\label{lem:row}
Let $G$ be a finite group and $\Lambda=\Z[G]$.  If $m\ge3$ and
\[
  x_1\Lambda+\cdots+x_m\Lambda=\Lambda,
\]
then finitely many elementary additions
\[
  x_j\longleftarrow x_j+x_iq,
  \qquad i\ne j,\quad q\in\Lambda,
\]
can make an entry of the row equal to $1$.
\end{lemma}

\begin{proof}
We first record how to correct a designated entry over a
finite-dimensional semisimple algebra $R$.  Suppose $xR+U=R$, where
$U$ is a right ideal.  Let $T=L_x:R\to R$ be left multiplication by
$x$, let $\pi:R\to R/U$ be the quotient map, and put
$K=\ker(\pi T)$.  The map $\pi T$ is onto, and the split exact
sequences give
\[
  R\cong K\oplus R/U,
  \qquad R\cong U\oplus R/U.
\]
Comparison of simple-module multiplicities therefore gives $K\cong U$.
Choose an isomorphism $h:K\to U$.  Extend
$h-T|_K:K\to U$ to a right-module map $c:R\to U$ using a complement
of $K$.  The map $T+c$ is an isomorphism: it restricts to $h$ on $K$
and induces the isomorphism $R/K\to R/U$ supplied by $\pi T$.
Every right-module endomorphism of $R$ is left multiplication by its
value at $1$.  Hence $T+c=L_{x+u}$ for $u=c(1)\in U$, and $x+u$ is
a unit.  If $U=\sum_{i\ne j}x_iR$, this corrects only the designated
entry $x_j$ by additions of right multiples of the other entries.

Apply this observation first to $R=\Q[G]$, which is semisimple by
Maschke's theorem, with designated entry $x_1$.  In group-basis
coordinates for the correcting coefficients $q_i$, the determinant
\[
  \det_{\Q} L_{x_1+\sum_{i\ne1}x_iq_i}
\]
is a polynomial with integer coefficients.  The unit correction over
$\Q[G]$ shows that it is nonzero.  A nonzero polynomial over $\Q$
does not vanish at every integer tuple, by induction on its number of
variables.  We may therefore choose all $q_i\in\Lambda$ with this
determinant nonzero.  After these elementary additions, rename the
first entry $x_1$.  Its principal right ideal has finite additive index
\[
  D=[\Lambda:x_1\Lambda]=|\det L_{x_1}|,
  \qquad D\Lambda\subset x_1\Lambda.
\]

We keep this entry fixed and correct only $x_2$.  For each prime
$p\mid D$, let $R_p=\Lambda/p\Lambda$ and
$J_p=\operatorname{Jac}(R_p)$.  The current row generates the unit
right ideal in the semisimple algebra $R_p/J_p$.  The preceding
observation gives coefficients $\bar q_{i,p}\in R_p$, $i\ne2$, such
that
\[
  \bar x_2+\sum_{i\ne2}\bar x_i\bar q_{i,p}
\]
is a unit modulo $J_p$, and hence a unit in $R_p$.  Here units lift
across the nilpotent ideal $J_p$, since $1-j$ has a finite geometric
series inverse for $j\in J_p$.  Choose $q_i\in\Lambda$, $i\ne2$,
with these residues for all $p\mid D$ by coefficientwise Chinese
remaindering, and make the additions
\[
  x_2\longleftarrow x_2+\sum_{i\ne2}x_iq_i.
\]
All source entries remain fixed during these additions, so the new
$x_2$ is a unit modulo every prime dividing the unchanged integer $D$.

The additive quotient
$A=\Lambda/(x_1\Lambda+x_2\Lambda)$ is finite and killed by $D$.
For each $p\mid D$, the invertibility of $x_2$ modulo $p$ gives
$x_2\Lambda+p\Lambda=\Lambda$, and hence $A=pA$.  Iteration gives
$A=DA=0$.  This also covers $D=1$, when $A$ is already zero.
We have proved $x_1\Lambda+x_2\Lambda=\Lambda$.

Finally put $t=x_3$ and choose $a,c\in\Lambda$ with $x_1a+x_2c=1$.  The two
elementary additions to the third entry with coefficients
$a(1-t)$ and $c(1-t)$ give
\[
  t+x_1a(1-t)+x_2c(1-t)=1.
\]
Neither source entry changes, and all products retain the required
order.
\end{proof}
